\documentclass[../main]{subfiles}
\newcounter{z}
\newcommand\y[2]{
  \loop \ifnum\value{z} < #1
    #2%
    \stepcounter{z}%
  \repeat
}

\begin{document}
\chapter{Pasantías}
\img{images/14/pasantias.jpg}{Alumnos realizando pasantías en Tecpetrol (Grupo Techint)}

\info{¿Qué es una pasantía?}{
Una pasantía estudiantil de secundaria técnica es una experiencia de aprendizaje práctica en la que los estudiantes de secundaria trabajan en empresas o entornos relacionados con su campo de estudio técnico o vocacional. Estas pasantías les permiten aplicar lo que han aprendido en el aula, obtener orientación de profesionales experimentados y adquirir habilidades prácticas, lo que enriquece su educación y les ayuda a tomar decisiones informadas sobre futuras carreras.}

\section{Pasantía}
Las pasantías ocupan un lugar sustancial en el currículum de la escuela técnica como estrategia de formación para los alumnos y como posibilidad de acercamiento al mundo del trabajo. 
A partir de un estudio de caso, basado en la obtención de datos primarios (entrevistas al equipo directivo, alumnos y empresarios) y al análisis de datos secundarios (documentos ministeriales e institucionales), se pudo identificar que, en el proceso de puesta en acto de la pasantía, la trayectoria escolar y el rendimiento académico de los alumnos actúan como mecanismos de selección para la realización de la pasantía en determinadas empresas. 

En consecuencia, la forma de acceso al sistema de pasantías genera prácticas diferenciadas e inequitativas en la formación de los alumnos. La empresa a la que asisten se constituye en un factor que influye en su formación y en sus expectativas respecto a la posibilidad de un rápido ingreso al mundo laboral.




\newpage

\subsection{Planilla de asistencia a pasantías}

\begin{table}[h]
    \centering
    \begin{tabular}{|c|c|c|p{5cm}|p{5cm}|} \hline 
         Fecha& Entrada &Salida & Firma Alumno& Comentarios  \\\hline \hline
       ..../..../2026 & : & : & & \\ \hline
       ..../..../2026 & : & : & & \\ \hline
       ..../..../2026 & : & : & & \\ \hline
       ..../..../2026 & : & : & & \\ \hline
       ..../..../2026 & : & : & & \\ \hline
       ..../..../2026 & : & : & & \\ \hline
       ..../..../2026 & : & : & & \\ \hline
       ..../..../2026 & : & : & & \\ \hline
       ..../..../2026 & : & : & & \\ \hline
       ..../..../2026 & : & : & & \\ \hline
       ..../..../2026 & : & : & & \\ \hline
       ..../..../2026 & : & : & & \\ \hline
       ..../..../2026 & : & : & & \\ \hline
       ..../..../2026 & : & : & & \\ \hline
       ..../..../2026 & : & : & & \\ \hline
       ..../..../2026 & : & : & & \\ \hline
       ..../..../2026 & : & : & & \\ \hline
       ..../..../2026 & : & : & & \\ \hline
       ..../..../2026 & : & : & & \\ \hline
       ..../..../2026 & : & : & & \\ \hline
       ..../..../2026 & : & : & & \\ \hline
       ..../..../2026 & : & : & & \\ \hline
       ..../..../2026 & : & : & & \\ \hline
       ..../..../2026 & : & : & & \\ \hline
       ..../..../2026 & : & : & & \\ \hline
       ..../..../2026 & : & : & & \\ \hline
       ..../..../2026 & : & : & & \\ \hline
       ..../..../2026 & : & : & & \\ \hline
       ..../..../2026 & : & : & & \\ \hline
 
      
    
    \end{tabular}
\end{table}
\vspace{2cm}
\begin{center}
\begin{tabular}{c c}
    ......................................... & .........................................    \\
            Firma Encargado/Jefe.  & Firma Docente \\
\end{tabular}
\end{center}
\newpage
\comment{}{
\subsection{Informe de pasantías.}

Las pasantías en el curso 5º3º son realizadas por 6 alumnos. Los alumnos realizan una pasantía de 4hs por día hábil, en el horario de tarde, hasta cumplimentar en una primera parte las 80hs de pasantías.

Las actividades que realizan están relacionadas con la industria y el comercio de los insumos de electricidad, donde elaboran tableros eléctricos interpretando circuitos y relacionándolos con componentes eléctricos del circuito comercial. Además, realizan el fraccionamiento de cables de baja tensión de diferentes tipos: unipolares, subterráneos y tripolares.

\begin{table}[h!]
    \centering
    \begin{tabular}{|p{7cm}|c|c|c|c|c|} \hline
        Apellido y Nombres de Alumnos  & Fecha inicio & Fecha fin & Lugar \\ \hline
         Alberti Gimenez, Gino Patricio & & &  Sierra Electricidad \\ \hline
         Alvarez, Lucas Nahuel & & & Sierra Electricidad\\ \hline
         Carretero Troncoso, Emiliano Fabian & & & Fase Electricidad\\ \hline
         Lomoro, Luciano Leonel & & & Sierra Electricidad\\ \hline
         Rodriguez Benitez, Juan Manuel & & & Fase Electricidad\\ \hline
         Sandoval, Juan Camilo &  & & Fase Electricidad\\ \hline 
         
      
    \end{tabular}
    \caption{Listado de pasantes.}
    \label{tab:my_label}
\end{table}
\newpage
\subsection{Fase electricidad}

\begin{table}[h!]
    \centering
    \begin{tabular}{|c|c|} \hline
        Razon social &  \\
        Domicilio & \\
        Telefono & \\
        Persona Contacto & \\
        Email & \\
        
    \end{tabular}
\end{table}

\subsection{Fase electricidad}
\subsubsection{Datos de la empresa.}
}





\end{document}